\documentclass[11pt,a4paper]{article}

% ==================== TEMEL PAKETLER (pdflatex) ====================
\usepackage[T1]{fontenc}
\usepackage[utf8]{inputenc}
\usepackage{lmodern}
\usepackage[a4paper,margin=2cm,top=2.4cm,bottom=2.4cm]{geometry}
\usepackage{xcolor}
\usepackage{colortbl}
\usepackage{array}
\usepackage{booktabs}
\usepackage{longtable}
\usepackage{tabularx}
\usepackage{listings}
\usepackage{fancyhdr}
\usepackage{amsmath}
\usepackage{amssymb}
\usepackage{hyperref}

% ==================== RENK PALETİ ====================
\definecolor{anaMavi}{HTML}{132C4A}
\definecolor{vurguMavi}{HTML}{2E6DA4}
\definecolor{acikMavi}{HTML}{EAF2F9}
\definecolor{turuncu}{HTML}{D9701E}
\definecolor{yesil}{HTML}{2E7D32}
\definecolor{acikYesil}{HTML}{E8F5E9}
\definecolor{kirmizi}{HTML}{C0392B}
\definecolor{mor}{HTML}{6C3483}
\definecolor{acikMor}{HTML}{F3E9F7}
\definecolor{acikGri}{HTML}{F4F5F7}
\definecolor{ortaGri}{HTML}{6B7280}
\definecolor{koyuGri}{HTML}{2D2D2D}
\definecolor{kodArka}{HTML}{FBFBFA}

% ==================== HYPERREF ====================
\hypersetup{
  colorlinks=true,
  linkcolor={[HTML]{2E6DA4}},
  urlcolor={[HTML]{2E6DA4}},
  pdftitle={Boost C++ Kütüphaneleri Kapsamlı Rehber},
  pdfauthor={C++ Kütüphaneleri Rehberi},
  bookmarksnumbered=true
}

% ==================== BAŞLIK STİLLERİ (titlesec olmadan) ====================
\setcounter{secnumdepth}{0}
\newcounter{secc}
\newcounter{subsecc}[secc]
\newcounter{subsubsecc}[subsecc]
\newcommand{\gsec}[1]{\par\phantomsection\stepcounter{secc}%
  \setcounter{subsecc}{0}%
  \vspace{18pt}{\color{anaMavi}\Large\bfseries\thesecc.\hspace{0.5em}#1}\par
  \vspace{2pt}{\color{turuncu}\rule{\linewidth}{1.6pt}}\par\vspace{8pt}%
  \addcontentsline{toc}{section}{\protect\numberline{\thesecc}#1}\markboth{#1}{#1}}
\newcommand{\gsub}[1]{\par\phantomsection\stepcounter{subsecc}%
  \setcounter{subsubsecc}{0}%
  \vspace{11pt}{\color{vurguMavi}\large\bfseries\thesecc.\thesubsecc\hspace{0.5em}#1}\par\vspace{4pt}%
  \addcontentsline{toc}{subsection}{\protect\numberline{\thesecc.\thesubsecc}#1}}
\newcommand{\gsubsub}[1]{\par\phantomsection\stepcounter{subsubsecc}%
  \vspace{7pt}{\color{koyuGri}\normalsize\bfseries #1}\par\vspace{2pt}}

% Kısım (Part) başlığı
\newcommand{\gpart}[2]{\clearpage
  \begin{center}\vspace*{1.5cm}
  {\color{turuncu}\rule{0.7\linewidth}{1.5pt}}\\[10pt]
  {\color{ortaGri}\Large KISIM #1}\\[8pt]
  {\color{anaMavi}\Huge\bfseries #2}\\[10pt]
  {\color{turuncu}\rule{0.7\linewidth}{1.5pt}}
  \end{center}\vspace{1cm}}

% ==================== BAŞLIK / ALTBİLGİ ====================
\setlength{\headheight}{14pt}
\pagestyle{fancy}
\fancyhf{}
\renewcommand{\headrulewidth}{0.4pt}
\renewcommand{\footrulewidth}{0.4pt}
\fancyhead[L]{\small\color{ortaGri}Boost C++ Kütüphaneleri}
\fancyhead[R]{\small\color{ortaGri}\leftmark}
\fancyfoot[C]{\small\color{ortaGri}\thepage}

% ==================== KOD LİSTELEME (C++) ====================
\lstdefinestyle{cppstyle}{
  language=C++,
  basicstyle=\ttfamily\small\color{koyuGri},
  keywordstyle=\color{vurguMavi}\bfseries,
  commentstyle=\color{yesil}\itshape,
  stringstyle=\color{turuncu},
  numberstyle=\tiny\color{ortaGri},
  numbers=left,
  numbersep=8pt,
  backgroundcolor=\color{kodArka},
  frame=single,
  framesep=6pt,
  rulecolor=\color{ortaGri!40},
  breaklines=true,
  showstringspaces=false,
  tabsize=4,
  columns=fullflexible,
  keepspaces=true,
  extendedchars=true,
  literate={ı}{{\i}}1{İ}{{\.I}}1{ş}{{\c s}}1{Ş}{{\c S}}1{ğ}{{\u g}}1{Ğ}{{\u G}}1{ç}{{\c c}}1{Ç}{{\c C}}1{ö}{{\"o}}1{Ö}{{\"O}}1{ü}{{\"u}}1{Ü}{{\"U}}1,
  morekeywords={size_t,uint8_t,uint16_t,uint32_t,uint64_t,int8_t,int16_t,
    int32_t,int64_t,nullptr,constexpr,noexcept,override,final,auto,
    vector,string,pair,map,set,unordered_map,unordered_set,priority_queue,
    queue,stack,deque,array,list,tuple,optional,shared_ptr,unique_ptr,
    template,typename,class,public,private,protected,virtual,namespace,
    using,std,boost,co_await,co_return,co_yield,concept,requires,
    move,make_pair,make_shared,make_unique,decltype}
}
\lstset{style=cppstyle}

% Kabuk/terminal stili
\lstdefinestyle{shstyle}{
  basicstyle=\ttfamily\small\color{koyuGri},
  backgroundcolor=\color{koyuGri!6},
  frame=single,framesep=6pt,rulecolor=\color{ortaGri!40},
  breaklines=true,showstringspaces=false,tabsize=4,columns=fullflexible,
  keepspaces=true,
  literate={ı}{{\i}}1{İ}{{\.I}}1{ş}{{\c s}}1{Ş}{{\c S}}1{ğ}{{\u g}}1{Ğ}{{\u G}}1{ç}{{\c c}}1{Ç}{{\c C}}1{ö}{{\"o}}1{Ö}{{\"O}}1{ü}{{\"u}}1{Ü}{{\"U}}1,
}

% ==================== ÖZEL KUTULAR (tcolorbox olmadan) ====================
\newcommand{\callout}[4]{\par\smallskip\noindent\begingroup
  \setlength{\fboxsep}{8pt}\setlength{\fboxrule}{0.8pt}%
  \fcolorbox{#1}{#2}{\parbox{\dimexpr\linewidth-2\fboxsep-2\fboxrule\relax}{%
    {\bfseries\color{#1}#3}\par\smallskip #4}}\endgroup\par\smallskip}
\newcommand{\bilgi}[1]{\callout{vurguMavi}{acikMavi}{Bilgi}{#1}}
\newcommand{\ipucu}[1]{\callout{yesil}{acikYesil}{İpucu}{#1}}
\newcommand{\uyari}[1]{\callout{kirmizi}{kirmizi!7}{Dikkat}{#1}}
\newcommand{\ornek}[2][Örnek]{\callout{ortaGri!90!black}{acikGri}{#1}{#2}}
\newcommand{\notk}[2][Not]{\callout{mor}{acikMor}{#1}{#2}}

% Yardımcı komutlar
\newcommand{\code}[1]{\texttt{\color{koyuGri}#1}}
\newcommand{\bigO}[1]{\ensuremath{\mathcal{O}(#1)}}
% Matematik sembollerini metin modunda da güvenli kıl
\let\svtimes\times   \renewcommand{\times}{\ensuremath{\svtimes}}
\let\svcdot\cdot     \renewcommand{\cdot}{\ensuremath{\svcdot}}
\let\svto\to         \renewcommand{\to}{\ensuremath{\svto}}
\let\svle\le         \renewcommand{\le}{\ensuremath{\svle}}
\let\svge\ge         \renewcommand{\ge}{\ensuremath{\svge}}
\renewcommand{\arraystretch}{1.35}

% ==================== BAŞLANGIÇ ====================
\begin{document}

% ---------- KAPAK ----------
\begin{titlepage}
\centering
\vspace*{1.4cm}
{\color{turuncu}\rule{\textwidth}{2pt}}\\[8pt]
{\color{anaMavi}\Huge\bfseries Boost C++ Kütüphaneleri}\\[10pt]
{\color{vurguMavi}\LARGE Kapsamlı Uygulamalı Rehber}\\[8pt]
{\color{ortaGri}\large Modern C++ için Endüstri Standardı Kütüphane Koleksiyonu}\\[6pt]
{\color{turuncu}\rule{\textwidth}{2pt}}\\[1.4cm]

\begin{center}
\fcolorbox{vurguMavi}{acikMavi}{\parbox{0.86\textwidth}{\centering\large
Bu belge, Boost kütüphanelerini modern C++ (C++17/20) ile,
her konuyu \textbf{derlenebilir kod örnekleriyle}, mimari açıklamalarla
ve pratik kullanım desenleriyle adım adım anlatır. Akıllı işaretçilerden
Boost.Asio ile ağ programlamaya, metaprogramlamadan graf algoritmalarına
kadar en çok kullanılan Boost bileşenlerini kapsar.\vspace{4pt}}}
\end{center}

\vspace{0.8cm}
\begin{center}
\fcolorbox{ortaGri!50}{acikGri}{\parbox{0.86\textwidth}{\small
\textbf{Kapsam --- Kısım 1 (Temeller):} Boost felsefesi \textbullet\ kurulum
\textbullet\ b2 \& CMake entegrasyonu \textbullet\ header-only vs derlenmiş.\vspace{4pt}

\textbf{Kapsam --- Kısım 2 (Bellek \& Yardımcılar):} akıllı işaretçiler
\textbullet\ Optional \textbullet\ Variant \textbullet\ Any \textbullet\ lexical\_cast
\textbullet\ Format \textbullet\ String Algo \textbullet\ Tokenizer.\vspace{4pt}

\textbf{Kapsam --- Kısım 3 (Konteynerler):} Container \textbullet\ MultiIndex
\textbullet\ Bimap \textbullet\ CircularBuffer \textbullet\ Intrusive
\textbullet\ DynamicBitset.\vspace{4pt}

\textbf{Kapsam --- Kısım 4 (Eşzamanlılık \& Ağ):} Thread \textbullet\ Atomic
\textbullet\ Lockfree \textbullet\ \textbf{Asio (detaylı)} \textbullet\ Fiber.\vspace{4pt}

\textbf{Kapsam --- Kısım 5 (Metin \& Veri):} Regex \textbullet\ Spirit
\textbullet\ PropertyTree \textbullet\ JSON \textbullet\ Serialization.\vspace{4pt}

\textbf{Kapsam --- Kısım 6 (Sistem, Matematik \& Graf):} Filesystem
\textbullet\ DateTime \textbullet\ ProgramOptions \textbullet\ Interprocess
\textbullet\ Math \textbullet\ Multiprecision \textbullet\ Graph (BGL).\vspace{4pt}}}
\end{center}

\vfill
{\color{ortaGri}\large Hazırlanma Tarihi: 19 Temmuz 2026}\\[4pt]
{\color{ortaGri} Boost 1.85+ \textbullet\ C++17 / C++20 \textbullet\ g++ / clang++}
\vspace*{0.6cm}
\end{titlepage}

% ---------- İÇİNDEKİLER ----------
\tableofcontents
\thispagestyle{empty}
\clearpage

% #####################################################################
\gpart{1}{Temeller, Felsefe ve Kurulum}
% #####################################################################

% =====================================================================
\gsec{Boost Nedir ve Neden Kullanılır?}

Boost, C++ dünyasının en etkili ve en saygın açık kaynak kütüphane
koleksiyonudur. 1998'de, C++ standart komitesinin (ISO WG21) bazı üyeleri
tarafından, standarda aday olabilecek yüksek kaliteli kütüphaneleri
geliştirmek ve olgunlaştırmak amacıyla kuruldu. Bugün 160'tan fazla
bağımsız kütüphaneden oluşan devasa bir ekosistemdir.

Boost'un C++ standardı üzerindeki etkisi tarihseldir: bugün standart
kütüphanenin (\code{std}) ayrılmaz parçaları olan pek çok bileşen önce
Boost'ta doğdu, yıllarca gerçek projelerde denendi ve olgunlaştıktan sonra
standarda alındı.

\begin{center}
\begin{tabular}{@{}>{\raggedright\arraybackslash}p{5.2cm} >{\raggedright\arraybackslash}p{5.2cm} c@{}}
\toprule
\textbf{Boost Kütüphanesi} & \textbf{Standart Karşılığı} & \textbf{Sürüm} \\
\midrule
\code{boost::shared\_ptr} & \code{std::shared\_ptr} & C++11 \\
\code{boost::function} & \code{std::function} & C++11 \\
\code{boost::bind} & \code{std::bind} / lambdalar & C++11 \\
\code{boost::array} & \code{std::array} & C++11 \\
\code{boost::thread} & \code{std::thread} & C++11 \\
\code{boost::optional} & \code{std::optional} & C++17 \\
\code{boost::variant} & \code{std::variant} & C++17 \\
\code{boost::any} & \code{std::any} & C++17 \\
\code{boost::filesystem} & \code{std::filesystem} & C++17 \\
\code{boost::string\_view} & \code{std::string\_view} & C++17 \\
\bottomrule
\end{tabular}
\end{center}

\bilgi{Boost'u öğrenmek, aslında C++'ın \emph{geleceğini} öğrenmektir.
Bugün Boost'ta kullandığınız birçok kütüphane, birkaç yıl sonra standardın
parçası olabilir. Ayrıca standardın henüz kapsamadığı devasa bir alan
(Asio, Spirit, Graph, Multiprecision, Geometry...) yalnızca Boost'ta mevcuttur.}

\gsub{Boost'u Ne Zaman Kullanmalı, Ne Zaman Kullanmamalı?}

\textbf{Kullanın:} Standart kütüphanenin karşılamadığı bir ihtiyaç varsa
(ağ, ayrıştırma, çok yüksek hassasiyetli aritmetik, graf algoritmaları,
tarih/saat işlemleri), yeni bir standart özelliği eski bir derleyicide
kullanmanız gerekiyorsa (C++17 \code{optional}'ı C++14 projesinde), veya
savaşta denenmiş, iyi test edilmiş bir çözüm istiyorsanız.

\textbf{Dikkatli olun:} Standart karşılığı zaten varsa ve C++17+ hedefliyorsanız,
önce \code{std} sürümünü tercih edin. Boost bağımlılığı derleme süresini
uzatır, ikili boyutu artırır ve dağıtımı zorlaştırabilir. Ayrıca bazı Boost
kütüphaneleri ağır şablon (template) metaprogramlama içerir ve derleme
hatalarını okumak zor olabilir.

\uyari{Boost \emph{tek bir kütüphane değildir}; birbirinden bağımsız
onlarca kütüphanenin şemsiye adıdır. ``Boost'u projeme ekledim'' demek
genellikle yanlıştır; ``Boost.Asio ve Boost.System'i ekledim'' demek doğrudur.
İhtiyacınız olan alt kütüphaneleri seçerek bağlayın.}

\gsub{Header-Only ve Derlenmiş Kütüphaneler}

Boost kütüphaneleri iki gruba ayrılır ve bu ayrım kurulum/derleme açısından
kritiktir:

\begin{itemize}
\item \textbf{Header-only (yalnızca başlık):} Tüm kodu başlık dosyalarında
bulunur. \code{\#include} etmeniz yeterlidir, ayrıca bağlamanız (link)
gerekmez. Çoğunluk buradadır: Optional, Variant, MPL, Hana, Spirit,
Lexical\_cast, Format, TypeTraits, Bind, Function, Signals2 (kısmen)...

\item \textbf{Derlenmiş (compiled/separately built):} Önceden derlenmiş
bir \code{.a}/\code{.so} kütüphaneye bağlanmayı gerektirir. Bunlar: System,
Filesystem, Thread, Regex, ProgramOptions, DateTime (kısmen), Serialization,
Log, Chrono, Context, Coroutine, IOStreams, Locale, Test (kısmen).
\end{itemize}

\ipucu{Bir kütüphanenin bağlanma gerektirip gerektirmediğini bilmiyorsanız,
Boost dokümantasyonundaki her kütüphanenin ``Getting Started'' sayfasına
bakın. Genel kural: dosya sistemi, iş parçacığı, ağ, günlükleme gibi
\emph{işletim sistemiyle etkileşen} kütüphaneler genelde derlenmiştir.}

% =====================================================================
\gsec{Kurulum ve Yapılandırma}

\gsub{Paket Yöneticisiyle Kurulum (En Kolay Yol)}

Çoğu Linux dağıtımında Boost paketlenmiştir. Geliştirme başlıkları için
\code{-dev} paketini kurmanız gerekir:

\begin{lstlisting}[style=shstyle]
# Debian / Ubuntu
sudo apt update
sudo apt install libboost-all-dev

# Fedora
sudo dnf install boost-devel

# Arch Linux / CachyOS
sudo pacman -S boost

# macOS (Homebrew)
brew install boost

# Sürümü doğrulama
dpkg -s libboost-dev | grep Version   # Debian ailesi
cat /usr/include/boost/version.hpp | grep BOOST_LIB_VERSION
\end{lstlisting}

\gsub{Kaynaktan Derleme (b2 / Boost.Build)}

En güncel sürümü veya özel yapılandırma istiyorsanız kaynaktan derleyin.
Boost, \code{b2} (eski adıyla \code{bjam}) adında kendi derleme sistemini
kullanır:

\begin{lstlisting}[style=shstyle]
# 1) Kaynağı indir ve aç
wget https://archives.boost.io/release/1.85.0/source/boost_1_85_0.tar.gz
tar xzf boost_1_85_0.tar.gz
cd boost_1_85_0

# 2) b2 aracını hazırla (bootstrap)
./bootstrap.sh --prefix=/usr/local --with-toolset=gcc

# 3) Derle ve kur (sadece ihtiyacınız olan kütüphaneleri seçebilirsiniz)
./b2 --with-system --with-filesystem --with-thread \
     variant=release link=shared,static \
     threading=multi cxxstd=17 -j$(nproc)

sudo ./b2 install       # başlıkları ve kütüphaneleri /usr/local'a kur

# TÜM kütüphaneleri derlemek isterseniz --with-* bayraklarını kaldırın
./b2 -j$(nproc) cxxstd=17
\end{lstlisting}

\begin{center}
\begin{tabular}{@{}ll@{}}
\toprule
\textbf{b2 seçeneği} & \textbf{Anlamı} \\
\midrule
\code{variant=release} & Optimizasyonlu (\code{-O2 -DNDEBUG}) sürüm \\
\code{variant=debug} & Hata ayıklama sembolleriyle \\
\code{link=static} & Statik kütüphane (\code{.a}) üret \\
\code{link=shared} & Paylaşımlı kütüphane (\code{.so}) üret \\
\code{threading=multi} & İş parçacığı güvenli sürüm \\
\code{cxxstd=17} & C++17 standardıyla derle \\
\code{-j\$(nproc)} & Tüm çekirdekleri kullan (paralel) \\
\code{--with-X} & Yalnızca X kütüphanesini derle \\
\code{--build-dir=/tmp/b} & Ara çıktı dizini \\
\bottomrule
\end{tabular}
\end{center}

\gsub{Elle Derleme: g++ Komut Satırı}

Header-only bir kütüphane kullanıyorsanız, tek yapmanız gereken başlık
yolunu vermektir:

\begin{lstlisting}[style=shstyle]
# Header-only (ör. Boost.Optional) -- bağlama YOK
g++ -std=c++17 -I/usr/local/include main.cpp -o program

# Derlenmiş kütüphane (ör. Filesystem + System) -- bağlama VAR
g++ -std=c++17 main.cpp -o program \
    -lboost_filesystem -lboost_system

# Özel kurulum yolu belirtme (kaynaktan derlediyseniz)
g++ -std=c++17 -I/usr/local/include main.cpp -o program \
    -L/usr/local/lib -lboost_filesystem -lboost_system \
    -Wl,-rpath,/usr/local/lib
\end{lstlisting}

\uyari{Bağlama sırası önemlidir: \code{-lboost\_filesystem} kütüphanesi
\code{-lboost\_system}'e bağımlıdır. GNU ld ile bağımlı kütüphaneyi
\emph{sonra} yazın: önce \code{filesystem}, sonra \code{system}. Yeni
Boost sürümlerinde Filesystem artık System'e derleme zamanı bağımlı
olmayabilir, ancak eski sürümlerde bu sıralama derleme hatalarını çözer.}

\gsub{CMake ile Boost Kullanımı (Önerilen Yöntem)}

Gerçek projelerde Boost'u CMake üzerinden bulmak en temiz yaklaşımdır.
\code{find\_package(Boost ...)} komutu bileşenleri bulur ve modern hedef
tabanlı (target-based) bağlama sağlar:

\begin{lstlisting}[style=shstyle]
cmake_minimum_required(VERSION 3.16)
project(BoostOrnek LANGUAGES CXX)

set(CMAKE_CXX_STANDARD 17)
set(CMAKE_CXX_STANDARD_REQUIRED ON)

# İhtiyacınız olan bileşenleri açıkça listeleyin.
# Header-only kütüphaneler COMPONENTS listesine YAZILMAZ.
find_package(Boost 1.75 REQUIRED COMPONENTS
    filesystem
    system
    program_options
)

add_executable(program main.cpp)

# Modern (target tabanlı) bağlama -- tercih edilen yol:
target_link_libraries(program PRIVATE
    Boost::filesystem
    Boost::system
    Boost::program_options
)

# Header-only bir kütüphane (ör. Boost.Optional) kullanıyorsanız,
# başlık yolları için genel Boost hedefini bağlayın:
target_link_libraries(program PRIVATE Boost::headers)
\end{lstlisting}

\ipucu{Boost'u CMake'in \code{FetchContent}'i ile de otomatik indirtip
derletebilirsiniz; ancak Boost çok büyük olduğu için bu, ilk derlemeyi
ciddi biçimde uzatır. Büyük projelerde sistem paketini veya önceden
derlenmiş bir Boost'u tercih edin. \code{Boost::headers} hedefi, yalnızca
başlık yolu gereken header-only kütüphaneler için idealdir.}

\gsub{İlk Boost Programı}

Boost'un doğru kurulduğunu doğrulayan minimal bir program yazalım.
Boost.Version ile sürümü, Boost.Lexical\_cast ile bir tür dönüşümünü test edelim:

\begin{lstlisting}
#include <boost/version.hpp>
#include <boost/lexical_cast.hpp>
#include <iostream>
#include <string>

int main() {
    // BOOST_VERSION: örn. 108500 => 1.85.0
    std::cout << "Boost surumu: "
              << BOOST_VERSION / 100000 << "."       // ana
              << BOOST_VERSION / 100 % 1000 << "."    // ara
              << BOOST_VERSION % 100 << "\n";         // yama

    // lexical_cast: string <-> sayi guvenli donusum
    int value = boost::lexical_cast<int>("42");
    std::string text = boost::lexical_cast<std::string>(3.14159);

    std::cout << "int:    " << value << "\n";
    std::cout << "string: " << text << "\n";
    return 0;
}
\end{lstlisting}

\begin{lstlisting}[style=shstyle]
# lexical_cast header-only oldugu icin bağlama gerekmez
g++ -std=c++17 first_program.cpp -o first_program
./first_program
# Cikti:
#   Boost surumu: 1.85.0
#   int:    42
#   string: 3.14159
\end{lstlisting}

% #####################################################################
\gpart{2}{Bellek Yönetimi ve Yardımcı Kütüphaneler}
% #####################################################################

% =====================================================================
\gsec{Akıllı İşaretçiler (Boost.SmartPtr)}

Boost.SmartPtr, C++'ta otomatik bellek yönetiminin öncüsüdür.
\code{std::shared\_ptr} ve \code{std::unique\_ptr} bu kütüphaneden doğdu.
Bugün C++11+ kullanıyorsanız çoğu durumda \code{std} sürümlerini tercih
edersiniz; ancak Boost hâlâ standartta \emph{olmayan} önemli akıllı
işaretçi türleri sunar.

\gsub{Standartta Bulunmayan Boost Akıllı İşaretçileri}

\begin{center}
\begin{tabular}{@{}>{\raggedright\arraybackslash}p{3.6cm} >{\raggedright\arraybackslash}p{8.4cm}@{}}
\toprule
\textbf{İşaretçi} & \textbf{Kullanım Amacı} \\
\midrule
\code{scoped\_ptr} & Tek sahiplik, taşınamaz (move edilemez). Basit RAII;
\code{unique\_ptr}'ın kısıtlı ama açık niyetli atası. \\
\code{intrusive\_ptr} & Referans sayacı \emph{nesnenin içinde} tutulur.
Bellek dostu; harici kontrol bloğu yoktur. \\
\code{local\_shared\_ptr} & Tek iş parçacıklı \code{shared\_ptr}; atomik
olmayan sayaç kullanır, daha hızlıdır. \\
\code{shared\_array} & Dizi için \code{shared\_ptr} (C++17 öncesi
\code{shared\_ptr<T[]>} yokken kritikti). \\
\bottomrule
\end{tabular}
\end{center}

\gsub{intrusive\_ptr: Gömülü Referans Sayacı}

\code{intrusive\_ptr}, referans sayacını kontrol nesnesinde değil, işaret
ettiği nesnenin \emph{kendi içinde} tutar. Bu, her nesne için ayrı bir
kontrol bloğu tahsis etmeyi önler; çok sayıda küçük nesne yönetiyorsanız
belleği ve önbellek yerelliğini iyileştirir. Karşılığında, nesnenizin
sayaç ve \code{intrusive\_ptr\_add\_ref} / \code{intrusive\_ptr\_release}
fonksiyonlarını sağlaması gerekir.

\begin{lstlisting}
#include <boost/intrusive_ptr.hpp>
#include <atomic>
#include <iostream>

// Referans sayaci taban sinifi -- nesneye "gomulu" sayac
class RefCounted {
public:
    RefCounted() : ref_count_(0) {}
    virtual ~RefCounted() = default;

    friend void intrusive_ptr_add_ref(RefCounted* p) noexcept {
        p->ref_count_.fetch_add(1, std::memory_order_relaxed);
    }
    friend void intrusive_ptr_release(RefCounted* p) noexcept {
        // release-acquire ile guvenli silme
        if (p->ref_count_.fetch_sub(1, std::memory_order_acq_rel) == 1) {
            delete p;   // son referans dustu
        }
    }
private:
    mutable std::atomic<int> ref_count_;
};

class Node : public RefCounted {
public:
    explicit Node(int v) : value(v) {
        std::cout << "Dugum(" << value << ") olusturuldu\n";
    }
    ~Node() override {
        std::cout << "Dugum(" << value << ") yok edildi\n";
    }
    int value;
};

int main() {
    // Sayac gomulu oldugu icin harici kontrol blogu YOK
    boost::intrusive_ptr<Node> a(new Node(10));
    {
        boost::intrusive_ptr<Node> b = a;   // sayac: 2
        std::cout << "b->deger = " << b->value << "\n";
    }   // b yok oldu -> sayac: 1 (nesne yasar)

    std::cout << "a hala gecerli: " << a->value << "\n";
    return 0;
}   // a yok oldu -> sayac: 0 -> "Dugum(10) yok edildi"
\end{lstlisting}

\bilgi{\code{intrusive\_ptr} özellikle oyun motorlarında, referans sayımlı
COM benzeri API'lerde ve milyonlarca küçük nesneli grafik/sahne
ağaçlarında tercih edilir. \code{shared\_ptr}'ın nesne başına ayrı kontrol
bloğu tahsisi bu senaryolarda ciddi ek yük getirir.}

\gsub{Boost.Pool: Havuz Tabanlı Bellek Ayırma}

Boost.Pool, aynı boyutta çok sayıda küçük nesneyi hızlıca tahsis/serbest
bırakmak için tasarlanmış bir bellek havuzudur. \code{new}/\code{delete}'in
her çağrıda işletim sistemine gitme maliyetini ortadan kaldırır.

\begin{lstlisting}
#include <boost/pool/object_pool.hpp>
#include <iostream>

struct Particle {
    float x, y, z;
    float speed;
    Particle(float px, float py, float pz)
        : x(px), y(py), z(pz), speed(0.f) {}
};

int main() {
    // object_pool: nesneleri havuzdan hizlica al/ver
    boost::object_pool<Particle> pool;

    // construct() bellegi havuzdan alir ve yapicIyi cagirir
    Particle* p1 = pool.construct(1.f, 2.f, 3.f);
    Particle* p2 = pool.construct(4.f, 5.f, 6.f);

    std::cout << "p1: " << p1->x << "," << p1->y << "\n";

    pool.destroy(p1);   // yikici cagrIlIr, bellek havuzda kalIr (yeniden kullanIlIr)
    // p2 icin destroy cagirmasak bile havuz yok olurken temizlenir
    return 0;
}   // havuz yok olurken kalan tum nesneler otomatik temizlenir
\end{lstlisting}

% =====================================================================
\gsec{Boost.Optional: Değeri Olmayabilen Nesneler}

\code{boost::optional<T>}, bir değerin ``var veya yok'' durumunu tip güvenli
biçimde ifade eder. C++17'de \code{std::optional} olarak standartlaştı, ancak
Boost sürümü daha eski derleyicilerde çalışır ve bazı ek özelliklere sahiptir
(örn. referansları sarmalayabilme).

\begin{lstlisting}
#include <boost/optional.hpp>
#include <boost/optional/optional_io.hpp>   // operator<< destegi
#include <iostream>
#include <string>
#include <map>

// Bir haritada anahtar bulunmayabilir -> optional ile ifade et
boost::optional<std::string> find(const std::map<int,std::string>& m, int key) {
    auto it = m.find(key);
    if (it == m.end())
        return boost::none;      // deger yok
    return it->second;           // deger var
}

int main() {
    std::map<int,std::string> dictionary = {{1,"bir"}, {2,"iki"}};

    boost::optional<std::string> r1 = find(dictionary, 1);
    boost::optional<std::string> r2 = find(dictionary, 99);

    if (r1)                                   // bool baglamI: deger var mI?
        std::cout << "Bulundu: " << *r1 << "\n";

    // value_or: yoksa varsayilan don
    std::cout << "99 -> " << r2.value_or("(yok)") << "\n";

    // get_value_or ve emplace
    boost::optional<int> number;
    std::cout << "Baslangic: " << number.get_value_or(-1) << "\n";
    number.emplace(42);                       // yerinde olustur
    std::cout << "Sonra: " << *number << "\n";

    // Boost'a ozgu: referans sarmalama (std::optional yapamaz)
    int x = 100;
    boost::optional<int&> ref = x;
    *ref = 200;                               // x'i degistirir
    std::cout << "x = " << x << "\n";         // 200
    return 0;
}
\end{lstlisting}

\bilgi{\code{boost::optional<T\&>} (referansa optional) standartta yoktur;
bu Boost'a özgü güçlü bir özelliktir. ``Belki bir nesneye başvuru'' ifade
etmek için ham işaretçiye (\code{T*}) alternatif, daha okunur bir yol sunar.}

\uyari{Bir \code{optional}'ı değeri yokken \code{*opt} veya \code{opt.value()}
ile açmak (dereference) tanımsız davranıştır (Boost) / istisna fırlatır
(std, \code{value()}). Her zaman önce \code{if (opt)} ile kontrol edin veya
\code{value\_or} kullanın.}

% =====================================================================
\gsec{Boost.Variant ve Variant2: Tip Güvenli Birleşimler}

\code{boost::variant}, birden fazla tipten \emph{tam olarak birini} tutabilen,
tip güvenli bir birleşim (union) sağlar. C++17'de \code{std::variant} oldu.
Boost ayrıca daha modern, istisna güvenli \code{boost::variant2::variant}
sunar (asla ``valueless'' duruma düşmez).

\begin{lstlisting}
#include <boost/variant.hpp>
#include <iostream>
#include <string>
#include <vector>

// Variant birden fazla tipten birini tutar
using Value = boost::variant<int, double, std::string>;

// Ziyaretci (visitor): her tip icin ayri davranis
struct PrintVisitor : public boost::static_visitor<void> {
    void operator()(int i) const {
        std::cout << "int: " << i << "\n";
    }
    void operator()(double d) const {
        std::cout << "double: " << d << "\n";
    }
    void operator()(const std::string& s) const {
        std::cout << "string: \"" << s << "\"\n";
    }
};

int main() {
    std::vector<Value> values = { 42, 3.14, std::string("merhaba") };

    // apply_visitor: aktif tipe gore dogru operator()'u cagirir
    for (const auto& d : values)
        boost::apply_visitor(PrintVisitor{}, d);

    // which(): aktif tipin indeksi (0-tabanli)
    Value d = std::string("test");
    std::cout << "Aktif indeks: " << d.which() << "\n";   // 2

    // get<T>: dogru tipse deger, degilse istisna
    try {
        int i = boost::get<int>(d);   // d string tutuyor -> hata
        (void)i;
    } catch (const boost::bad_get& e) {
        std::cout << "Yanlis tip: " << e.what() << "\n";
    }
    return 0;
}
\end{lstlisting}

\ipucu{Variant, \emph{toplam tipler} (sum types) veya \emph{ayrık birleşimler}
(discriminated unions) için idealdir: bir JSON değeri (sayı/metin/dizi/nesne),
bir AST düğümü, veya bir durum makinesinin durumları. Ziyaretçi deseni,
her yeni tip eklendiğinde derleyicinin sizi eksik durumları ele almaya
zorlaması sayesinde çok güvenlidir.}

\notk[Modern Alternatif]{Yeni kod yazıyorsanız ve istisna güvenliği
kritikse, \code{boost::variant2} veya \code{std::variant} tercih edin.
Klasik \code{boost::variant}, bir atama sırasında istisna olursa nadiren
``valueless by exception'' duruma düşebilir; \code{variant2} bunu
tasarımıyla engeller.}

% =====================================================================
\gsec{Boost.Any: Herhangi Bir Tipi Tutmak}

\code{boost::any}, \emph{herhangi} bir tipteki tek bir değeri tip-silinmiş
(type-erased) biçimde tutar. Variant'tan farkı: önceden tip listesi
belirtmezsiniz. C++17'de \code{std::any} oldu.

\begin{lstlisting}
#include <boost/any.hpp>
#include <iostream>
#include <string>
#include <vector>

int main() {
    // Ayni kapta farkli tipler (heterojen)
    std::vector<boost::any> container;
    container.push_back(42);
    container.push_back(std::string("metin"));
    container.push_back(3.14);

    for (const auto& a : container) {
        // any_cast<T>: dogru tipse deger, degilse bad_any_cast
        if (a.type() == typeid(int))
            std::cout << "int: " << boost::any_cast<int>(a) << "\n";
        else if (a.type() == typeid(std::string))
            std::cout << "string: " << boost::any_cast<std::string>(a) << "\n";
        else if (a.type() == typeid(double))
            std::cout << "double: " << boost::any_cast<double>(a) << "\n";
    }

    // Isaretci ile guvenli cast: yanlissa nullptr (istisna atmaz)
    boost::any a = 100;
    if (int* p = boost::any_cast<int>(&a))
        std::cout << "Isaretci cast: " << *p << "\n";
    return 0;
}
\end{lstlisting}

\uyari{\code{any}, tip güvenliğini çalışma zamanına erteler ve her erişimde
tip kontrolü gerektirir; ayrıca küçük olmayan nesneler için yığın (heap)
tahsisi yapabilir. Tipler önceden biliniyorsa \code{variant} neredeyse her
zaman daha hızlı ve güvenlidir. \code{any}'yi yalnızca gerçekten
\emph{açık uçlu} tip kümeleri için kullanın (örn. eklenti sistemleri).}

% =====================================================================
\gsec{Metin İşleme Yardımcıları}

\gsub{Boost.Lexical\_cast: Basit Tip Dönüşümleri}

\code{lexical\_cast}, metin ile sayısal tipler arasında \code{<<}/\code{>>}
akış operatörlerini kullanarak dönüşüm yapar. \code{std::to\_string} veya
\code{std::stoi}'den daha genel; herhangi bir akış-uyumlu tip için çalışır.

\begin{lstlisting}
#include <boost/lexical_cast.hpp>
#include <iostream>
#include <string>

int main() {
    // string -> sayi
    int    i = boost::lexical_cast<int>("123");
    double d = boost::lexical_cast<double>("3.14");

    // sayi -> string
    std::string s = boost::lexical_cast<std::string>(2026);

    std::cout << i << " " << d << " " << s << "\n";

    // Hatali girdi -> bad_lexical_cast istisnasi
    try {
        int error = boost::lexical_cast<int>("abc");
        (void)error;
    } catch (const boost::bad_lexical_cast& e) {
        std::cout << "Donusum hatasi: " << e.what() << "\n";
    }
    return 0;
}
\end{lstlisting}

\gsub{Boost.Format: Tip Güvenli printf}

\code{boost::format}, \code{printf}'in tip güvenli ve genişletilebilir
alternatifidir. \code{\%} operatörüyle argümanları besler; yanlış tip
sorunlarını ve bellek taşmalarını önler.

\begin{lstlisting}
#include <boost/format.hpp>
#include <iostream>
#include <string>

int main() {
    // %1%, %2% ... konumsal argumanlar
    std::string s = (boost::format("%1% + %2% = %3%") % 2 % 3 % 5).str();
    std::cout << s << "\n";                       // 2 + 3 = 5

    // printf tarzi bicimlendirme belirteclerI
    std::cout << boost::format("Pi ~ %.4f\n") % 3.14159265;
    std::cout << boost::format("Hex: %1$#x, Onluk: %1$d\n") % 255;

    // Alan genisligi ve hizalama
    std::cout << boost::format("|%1$10s|%2$-10s|\n") % "sag" % "sol";

    // Ayni argumani tekrar kullanma (%1% birden fazla kez)
    std::cout << boost::format("%1%^2 = %2%\n") % 7 % (7*7);
    return 0;
}
\end{lstlisting}

\ipucu{\code{boost::format}, konumsal argümanları (\code{\%1\%}) desteklediği
için \emph{yerelleştirmede} (i18n) çok kullanışlıdır: farklı dillerde
kelime sırası değiştiğinde biçim dizesini değiştirmeniz yeterlidir,
kod aynı kalır.}

\gsub{Boost.String\_algo: Metin Algoritmaları}

Boost.String\_algo, standart kütüphanede eksik olan onlarca metin işleme
fonksiyonu sunar: büyük/küçük harf, kırpma (trim), bölme (split),
birleştirme (join), değiştirme (replace), arama.

\begin{lstlisting}
#include <boost/algorithm/string.hpp>
#include <iostream>
#include <string>
#include <vector>

int main() {
    namespace ba = boost::algorithm;

    std::string s = "  Merhaba, Dunya  ";

    // Kirpma (bosluklari kaldir) -- yerinde (in-place)
    std::string t = s;
    ba::trim(t);
    std::cout << "[" << t << "]\n";               // [Merhaba, Dunya]

    // Buyuk/kucuk harf
    std::cout << ba::to_upper_copy(t) << "\n";    // MERHABA, DUNYA
    std::cout << ba::to_lower_copy(t) << "\n";    // merhaba, dunya

    // Bolme (split) -> vektore
    std::vector<std::string> parts;
    std::string csv = "elma,armut,kiraz,uzum";
    ba::split(parts, csv, ba::is_any_of(","));
    for (const auto& p : parts) std::cout << "- " << p << "\n";

    // Birlestirme (join)
    std::cout << ba::join(parts, " | ") << "\n";

    // Degistirme
    std::string text = "kedi kedi kedi";
    ba::replace_all(text, "kedi", "kopek");
    std::cout << text << "\n";                      // kopek kopek kopek

    // Baslangic/bitis/icerik kontrolu (buyuk-kucuk duyarsIz)
    std::cout << std::boolalpha
              << ba::istarts_with("Merhaba", "MER") << " "
              << ba::contains("Merhaba", "rha") << "\n";
    return 0;
}
\end{lstlisting}

\gsub{Boost.Tokenizer: Esnek Belirteçleme}

\code{boost::tokenizer}, bir dizeyi belirteçlere (token) ayırmanın esnek
bir yolunu sunar; farklı ayırıcı stratejileri (karakter, escape'li, CSV)
takılabilir.

\begin{lstlisting}
#include <boost/tokenizer.hpp>
#include <iostream>
#include <string>

int main() {
    std::string s = "bir,iki;uc dort";

    // Birden fazla ayirici karakter
    boost::char_separator<char> separator(",; ");
    boost::tokenizer<boost::char_separator<char>> tok(s, separator);

    for (const auto& t : tok)
        std::cout << "[" << t << "]";
    std::cout << "\n";        // [bir][iki][uc][dort]

    // CSV ayirici: tirnak ve escape'leri dogru isler
    std::string csv = "\"Ali, Veli\",42,\"istanbul\"";
    boost::tokenizer<boost::escaped_list_separator<char>> csvtok(csv);
    for (const auto& field : csvtok)
        std::cout << "Alan: " << field << "\n";
    return 0;
}
\end{lstlisting}

% #####################################################################
\gpart{3}{Konteynerler ve Veri Yapıları}
% #####################################################################

% =====================================================================
\gsec{Boost.Container: Gelişmiş Konteynerler}

Boost.Container, standart konteynerlerin geliştirilmiş sürümlerini ve
standartta bulunmayan yeni konteynerler sunar. Öne çıkanlar:

\begin{center}
\begin{tabular}{@{}>{\raggedright\arraybackslash}p{3.8cm} >{\raggedright\arraybackslash}p{8.2cm}@{}}
\toprule
\textbf{Konteyner} & \textbf{Özelliği} \\
\midrule
\code{flat\_map} / \code{flat\_set} & Sıralı vektör tabanlı; daha iyi önbellek
yerelliği, daha az bellek. Arama \bigO{\log n}, ekleme \bigO{n}. \\
\code{small\_vector} & Küçük boyutlar için yığın (stack) üzerinde depolama;
tahsisi geciktirir. \\
\code{static\_vector} & Sabit kapasiteli, hiç yığın tahsisi yapmayan vektör. \\
\code{stable\_vector} & Elemanlara işaretçi/iteratör kararlılığı garanti eder. \\
\code{devector} & Çift uçtan büyüyebilen vektör (deque + vector melezi). \\
\bottomrule
\end{tabular}
\end{center}

\begin{lstlisting}
#include <boost/container/flat_map.hpp>
#include <boost/container/small_vector.hpp>
#include <boost/container/static_vector.hpp>
#include <iostream>

int main() {
    namespace bc = boost::container;

    // flat_map: sirali vektor tabanli map. Az eleman + cok arama icin ideal.
    bc::flat_map<int, std::string> fm;
    fm[3] = "uc"; fm[1] = "bir"; fm[2] = "iki";
    for (const auto& [k, v] : fm)          // otomatik sirali: 1,2,3
        std::cout << k << "=" << v << " ";
    std::cout << "\n";

    // small_vector<T, N>: ilk N eleman stack'te, tasma olursa heap'e gecer
    bc::small_vector<int, 4> sv;
    sv.push_back(1); sv.push_back(2);      // heap tahsisi YOK (4'e kadar)
    std::cout << "small_vector boyut: " << sv.size() << "\n";

    // static_vector<T, N>: kapasite sabit N, ASLA heap kullanmaz
    bc::static_vector<int, 3> stv;
    stv.push_back(10); stv.push_back(20); stv.push_back(30);
    // stv.push_back(40);  // -> kapasite asIldI, tanImsIz/istisna
    std::cout << "static_vector kapasite: " << stv.capacity() << "\n";
    return 0;
}
\end{lstlisting}

\bilgi{\code{flat\_map}, \code{std::map}'e göre çok daha iyi önbellek
davranışı sergiler çünkü elemanlar bellekte bitişiktir. Okuma ağırlıklı,
nadiren değişen sözlükler için idealdir. Ancak sık ekleme/silme yaparsanız
her işlem \bigO{n} eleman kaydırdığından \code{std::map} daha iyidir.
C++23'te \code{std::flat\_map} olarak standarda girmiştir.}

% =====================================================================
\gsec{Boost.MultiIndex: Çok İndeksli Konteynerler}

Boost.MultiIndex, tek bir eleman kümesine \emph{aynı anda birden fazla
farklı indeksle} erişmenizi sağlar; tıpkı bir veritabanı tablosunun birden
fazla indeksi olması gibi. Bir çalışan listesini hem ID'ye, hem isme, hem
de maaşa göre sıralı tutmak isterseniz, veriyi kopyalamadan tek yapıda
yönetirsiniz.

\begin{lstlisting}
#include <boost/multi_index_container.hpp>
#include <boost/multi_index/ordered_index.hpp>
#include <boost/multi_index/member.hpp>
#include <iostream>
#include <string>

struct Employee {
    int id;
    std::string name;
    double salary;
    Employee(int i, std::string n, double m)
        : id(i), name(std::move(n)), salary(m) {}
};

// Indeksleri isimlendirmek icin etiketler
struct byId {};
struct byName {};
struct bySalary {};

namespace bmi = boost::multi_index;

// Uc farkli indeksli konteyner tanImI
using EmployeeTable = bmi::multi_index_container<
    Employee,
    bmi::indexed_by<
        // 1) ID'ye gore essiz (unique) sirali indeks
        bmi::ordered_unique<
            bmi::tag<byId>,
            bmi::member<Employee, int, &Employee::id>>,
        // 2) Isme gore (tekrar edebilir) sirali indeks
        bmi::ordered_non_unique<
            bmi::tag<byName>,
            bmi::member<Employee, std::string, &Employee::name>>,
        // 3) Maasa gore sirali indeks
        bmi::ordered_non_unique<
            bmi::tag<bySalary>,
            bmi::member<Employee, double, &Employee::salary>>
    >
>;

int main() {
    EmployeeTable table;
    table.insert(Employee(3, "Ayse",  8500));
    table.insert(Employee(1, "Mehmet", 7200));
    table.insert(Employee(2, "Zeynep", 9100));

    // ID'ye gore sirali gezinme
    std::cout << "== ID sirasi ==\n";
    for (const auto& c : table.get<byId>())
        std::cout << c.id << " " << c.name << " " << c.salary << "\n";

    // Isme gore arama
    auto& nameIdx = table.get<byName>();
    auto it = nameIdx.find("Zeynep");
    if (it != nameIdx.end())
        std::cout << "\nZeynep'in maasi: " << it->salary << "\n";

    // Maasa gore sirali (en dusukten en yuksege)
    std::cout << "\n== Maas sirasi ==\n";
    for (const auto& c : table.get<bySalary>())
        std::cout << c.salary << " -> " << c.name << "\n";
    return 0;
}
\end{lstlisting}

\ipucu{MultiIndex, bellekteki verileri bir mini-veritabanı gibi yönetmek
istediğinizde muhteşemdir: LRU önbellekleri, oyun varlık yöneticileri,
ağ bağlantı tabloları. \code{ordered}, \code{hashed}, \code{sequenced}
ve \code{random\_access} indeks türlerini karıştırabilirsiniz. Tek bir
elemanı güncellemek tüm indeksleri otomatik tutarlı tutar.}

% =====================================================================
\gsec{Boost.Bimap: Çift Yönlü Haritalar}

\code{boost::bimap}, hem anahtardan değere hem de değerden anahtara
\bigO{\log n} erişim sağlayan çift yönlü bir haritadır. Normal
\code{std::map} yalnızca tek yönde arama yaparken, bimap her iki yönü de
indeksler.

\begin{lstlisting}
#include <boost/bimap.hpp>
#include <iostream>
#include <string>

int main() {
    // Ulke <-> baskent iki yonlu esleme
    using Bimap = boost::bimap<std::string, std::string>;
    Bimap b;

    b.insert({"Turkiye", "Ankara"});
    b.insert({"Fransa",  "Paris"});
    b.insert({"Japonya", "Tokyo"});

    // Sol -> sag (ulke -> baskent)
    std::cout << "Turkiye'nin baskenti: "
              << b.left.at("Turkiye") << "\n";

    // Sag -> sol (baskent -> ulke)
    std::cout << "Paris hangi ulkede: "
              << b.right.at("Paris") << "\n";

    // Tum eslemeleri gezme
    for (const auto& c : b.left)
        std::cout << c.first << " <-> " << c.second << "\n";
    return 0;
}
\end{lstlisting}

\bilgi{Bimap, iki değer arasında birebir (one-to-one) ilişki olduğunda
idealdir: kullanıcı adı $\leftrightarrow$ ID, sembol $\leftrightarrow$ kod,
enum $\leftrightarrow$ metin. İki ayrı \code{std::map} tutup senkron kalmaya
çalışmaktan çok daha güvenlidir.}

% =====================================================================
\gsec{Boost.CircularBuffer: Dairesel Tampon}

\code{boost::circular\_buffer}, sabit kapasiteli, dolduğunda en eski
elemanların üzerine yazan bir tampondur. Sabit boyutlu kayan pencereler,
son N olayı tutma, ses/sinyal tamponları için idealdir.

\begin{lstlisting}
#include <boost/circular_buffer.hpp>
#include <iostream>

int main() {
    // Kapasite 3'lük dairesel tampon
    boost::circular_buffer<int> cb(3);

    cb.push_back(1);
    cb.push_back(2);
    cb.push_back(3);
    std::cout << "Dolu mu? " << std::boolalpha << cb.full() << "\n"; // true

    // Kapasite dolu; yeni eleman EN ESKIyi (1) atar
    cb.push_back(4);        // [2, 3, 4]
    cb.push_back(5);        // [3, 4, 5]

    std::cout << "Icerik: ";
    for (int x : cb) std::cout << x << " ";      // 3 4 5
    std::cout << "\n";

    std::cout << "En eski: " << cb.front()       // 3
              << ", en yeni: " << cb.back() << "\n";  // 5
    return 0;
}
\end{lstlisting}

% =====================================================================
\gsec{Boost.DynamicBitset: Dinamik Bit Kümeleri}

\code{std::bitset} derleme zamanında sabit boyutludur; \code{boost::dynamic\_bitset}
ise çalışma zamanında boyutlandırılabilir. Bit maskeleri, kümeler, elek
algoritmaları için idealdir.

\begin{lstlisting}
#include <boost/dynamic_bitset.hpp>
#include <iostream>

int main() {
    // Calisma zamanInda 16 bitlik kume
    boost::dynamic_bitset<> bits(16);

    bits.set(2);          // 2. biti 1 yap
    bits.set(5);
    bits.set(8);

    std::cout << "Bitler: " << bits << "\n";        // ...0100100100
    std::cout << "Kac bit 1: " << bits.count() << "\n";  // 3
    std::cout << "Boyut: " << bits.size() << "\n";       // 16

    // Bitsel operatorler
    boost::dynamic_bitset<> mask(16);
    mask.set(2); mask.set(3);
    auto intersection = bits & mask;    // AND
    std::cout << "Kesisim sayisi: " << intersection.count() << "\n";  // 1 (bit 2)

    // Ilk 1 olan biti bul, sonrakileri gez
    for (size_t i = bits.find_first();
         i != boost::dynamic_bitset<>::npos;
         i = bits.find_next(i)) {
        std::cout << "Set bit @ " << i << "\n";
    }
    return 0;
}
\end{lstlisting}

% #####################################################################
\gpart{4}{Eşzamanlılık ve Ağ Programlama}
% #####################################################################

% =====================================================================
\gsec{Boost.Thread: İş Parçacığı Yönetimi}

Boost.Thread, \code{std::thread}'in atasıdır ve standartta \emph{hâlâ
bulunmayan} bazı özellikler sunar: \code{shared\_mutex}'in daha zengin
sürümleri, kesintiye uğratılabilir (interruptible) iş parçacıkları,
\code{thread\_group}, ve gelişmiş \code{future}/\code{promise} desteği
(örn. \code{when\_all}, \code{then} zincirlemesi).

\begin{lstlisting}
#include <boost/thread.hpp>
#include <boost/thread/future.hpp>
#include <iostream>
#include <vector>

int compute(int x) {
    boost::this_thread::sleep_for(boost::chrono::milliseconds(50));
    return x * x;
}

int main() {
    // 1) Temel is parcacigi
    boost::thread t([]{
        std::cout << "Is parcacigindan merhaba\n";
    });
    t.join();

    // 2) thread_group: birden fazla is parcacigini birlikte yonet
    boost::thread_group group;
    for (int i = 0; i < 4; ++i)
        group.create_thread([i]{
            std::cout << "Isci " << i << " calisiyor\n";
        });
    group.join_all();     // hepsini bekle

    // 3) future/promise ile asenkron sonuc
    boost::future<int> f = boost::async(boost::launch::async, compute, 7);
    std::cout << "7^2 = " << f.get() << "\n";      // 49

    // 4) .then() ile surekli zincirleme (Boost'a ozgu, std'de yok)
    boost::future<int> chain =
        boost::async([]{ return 10; })
            .then([](boost::future<int> prev){
                return prev.get() * 2;             // 20
            });
    std::cout << "Zincir sonucu: " << chain.get() << "\n";
    return 0;
}
\end{lstlisting}

\uyari{Boost.Thread'in \code{.then()}, \code{when\_all}, \code{when\_any}
sürekli zincirleme (continuation) özellikleri için \code{BOOST\_THREAD\_PROVIDES\_FUTURE\_CONTINUATION}
gibi makrolar veya doğru derleme bayrakları gerekebilir. std::future bu
özelliklere sahip değildir; bu yüzden asenkron zincirleme istiyorsanız
Boost.Thread hâlâ değerlidir.}

\gsub{Karşılıklı Dışlama ve Kilitler}

\begin{lstlisting}
#include <boost/thread.hpp>
#include <iostream>

class SafeCounter {
    mutable boost::shared_mutex mtx_;   // okuyucu-yazici kilidi
    int value_ = 0;
public:
    void increment() {
        // Yazma: ozel (exclusive) kilit
        boost::unique_lock<boost::shared_mutex> lk(mtx_);
        ++value_;
    }
    int read() const {
        // Okuma: paylasimli (shared) kilit -- coklu okuyucu ayni anda
        boost::shared_lock<boost::shared_mutex> lk(mtx_);
        return value_;
    }
};

int main() {
    SafeCounter counter;
    boost::thread_group group;
    for (int i = 0; i < 10; ++i)
        group.create_thread([&]{
            for (int j = 0; j < 1000; ++j) counter.increment();
        });
    group.join_all();
    std::cout << "Toplam: " << counter.read() << "\n";   // 10000
    return 0;
}
\end{lstlisting}

% =====================================================================
\gsec{Boost.Asio: Asenkron Giriş/Çıkış ve Ağ}

Boost.Asio, C++'ta ağ ve düşük seviye asenkron G/Ç programlamanın
\emph{fiili standardıdır}. TCP/UDP soketleri, zamanlayıcılar, seri portlar,
ve platformdan bağımsız bir olay döngüsü (event loop) sunar. C++
Networking TS'nin temelidir. Asio'nun merkezinde \code{io\_context} (olay
döngüsü) bulunur; tüm asenkron işlemler ona kayıtlanır.

\gsub{Temel Kavramlar: io\_context ve Zamanlayıcılar}

\begin{lstlisting}
#include <boost/asio.hpp>
#include <iostream>

int main() {
    boost::asio::io_context io;

    // Asenkron zamanlayici: 2 saniye sonra tetiklenir
    boost::asio::steady_timer timer(io, boost::asio::chrono::seconds(2));

    timer.async_wait([](const boost::system::error_code& ec){
        if (!ec) std::cout << "2 saniye gecti!\n";
    });

    std::cout << "Bekleniyor...\n";
    io.run();   // olay dongusu: kayitli tum isler bitene kadar calisir
    std::cout << "io_context bitti\n";
    return 0;
}
\end{lstlisting}

\bilgi{\code{io\_context.run()} çağrısı, kayıtlı asenkron iş kalmayana
kadar bloke olur ve bu sırada tamamlanan işlerin geri çağırma (callback)
fonksiyonlarını çalıştırır. Asio'nun gücü, \emph{tek bir iş parçacığında}
binlerce eşzamanlı bağlantıyı yönetebilmesidir (proaktör deseni). Ağır yük
için \code{io.run()}'ı birden fazla iş parçacığında çağırarak ölçekleyebilirsiniz.}

\gsub{Asenkron TCP Yankı (Echo) Sunucusu}

Aşağıda, gelen her bağlantıyı kabul edip aldığı veriyi geri yollayan
tam bir asenkron TCP sunucusu bulunmaktadır. Bu, Asio'nun gerçek gücünü
gösteren kanonik örnektir.

\begin{lstlisting}
#include <boost/asio.hpp>
#include <iostream>
#include <memory>

using boost::asio::ip::tcp;

// Her baglanti icin bir oturum. shared_ptr ile omru yonetilir:
// asenkron isler tamamlanana kadar nesne yasar (enable_shared_from_this).
class Session : public std::enable_shared_from_this<Session> {
public:
    explicit Session(tcp::socket socket) : socket_(std::move(socket)) {}

    void start() { read(); }

private:
    void read() {
        auto self = shared_from_this();   // omru uzat
        socket_.async_read_some(
            boost::asio::buffer(buffer_),
            [this, self](boost::system::error_code ec, std::size_t n) {
                if (!ec) {
                    write(n);             // okunan veriyi geri yaz
                }
                // ec varsa oturum sessizce sonlanir (self referansi dusunce silinir)
            });
    }

    void write(std::size_t length) {
        auto self = shared_from_this();
        boost::asio::async_write(
            socket_,
            boost::asio::buffer(buffer_, length),
            [this, self](boost::system::error_code ec, std::size_t) {
                if (!ec) read();          // yazma bitti, tekrar oku (dongu)
            });
    }

    tcp::socket socket_;
    char buffer_[1024];
};

// Yeni baglantIlari kabul eden sunucu
class Server {
public:
    Server(boost::asio::io_context& io, unsigned short port)
        : acceptor_(io, tcp::endpoint(tcp::v4(), port)) {
        doAccept();
    }

private:
    void doAccept() {
        acceptor_.async_accept(
            [this](boost::system::error_code ec, tcp::socket socket) {
                if (!ec) {
                    std::cout << "Yeni baglanti: "
                              << socket.remote_endpoint() << "\n";
                    std::make_shared<Session>(std::move(socket))->start();
                }
                doAccept();  // bir sonraki baglantIyI kabul etmeye devam et
            });
    }
    tcp::acceptor acceptor_;
};

int main() {
    try {
        boost::asio::io_context io;
        Server server(io, 12345);
        std::cout << "Yanki sunucusu 12345 portunda dinliyor...\n";
        io.run();
    } catch (const std::exception& e) {
        std::cerr << "Hata: " << e.what() << "\n";
    }
    return 0;
}
\end{lstlisting}

\begin{lstlisting}[style=shstyle]
# Derleme (Asio, System kutuphanesine baglanir)
g++ -std=c++17 echo_server.cpp -o echo_server -lboost_system -lpthread
./echo_server &
# Baska bir terminalde test:
echo "merhaba asio" | nc localhost 12345
# Cikti: merhaba asio  (sunucu geri yansitir)
\end{lstlisting}

\uyari{Asio ile asenkron programlamada \code{shared\_from\_this} deseni
kritiktir: bir asenkron işlem devam ederken nesnenin (oturumun) yok
edilmemesini garanti eder. Geri çağırma lambdasına \code{self} kopyalayarak
referans sayacını yüksek tutarsınız; işlem bitince \code{self} kapsam
dışına çıkar ve nesne uygun şekilde temizlenir. Bu deseni atlarsanız
sarkan işaretçi (dangling pointer) ve çökme yaşarsınız.}

\gsub{Asenkron TCP İstemcisi}

\begin{lstlisting}
#include <boost/asio.hpp>
#include <iostream>
#include <string>

using boost::asio::ip::tcp;

int main() {
    try {
        boost::asio::io_context io;

        // Sunucu adresini coz (resolve)
        tcp::resolver resolver(io);
        auto endpoints = resolver.resolve("localhost", "12345");

        // Baglan (senkron -- basitlik icin)
        tcp::socket socket(io);
        boost::asio::connect(socket, endpoints);

        // Veri gonder
        std::string message = "Merhaba sunucu!";
        boost::asio::write(socket, boost::asio::buffer(message));

        // Yaniti oku
        char response[128];
        std::size_t n = socket.read_some(boost::asio::buffer(response));
        std::cout << "Sunucudan: " << std::string(response, n) << "\n";
    } catch (const std::exception& e) {
        std::cerr << "Istemci hatasi: " << e.what() << "\n";
    }
    return 0;
}
\end{lstlisting}

\gsub{Modern Yaklaşım: Coroutine'lerle Asio (C++20)}

Asio, C++20 coroutine'leriyle birlikte kullanıldığında asenkron kodu
\emph{senkron kod kadar okunur} hale getirir. \code{co\_await} ile geri
çağırma cehennemi (callback hell) tamamen ortadan kalkar.

\begin{lstlisting}
#include <boost/asio.hpp>
#include <boost/asio/co_spawn.hpp>
#include <boost/asio/detached.hpp>
#include <iostream>

using boost::asio::ip::tcp;
using boost::asio::awaitable;
using boost::asio::use_awaitable;
namespace this_coro = boost::asio::this_coro;

// Coroutine tabanli yanki oturumu -- geri caGirma yok, duz akIs!
awaitable<void> echo(tcp::socket socket) {
    char data[1024];
    try {
        for (;;) {
            // co_await: asenkron ama senkron gorunumlu
            std::size_t n = co_await socket.async_read_some(
                boost::asio::buffer(data), use_awaitable);
            co_await boost::asio::async_write(
                socket, boost::asio::buffer(data, n), use_awaitable);
        }
    } catch (const std::exception& e) {
        std::cout << "Oturum bitti: " << e.what() << "\n";
    }
}

awaitable<void> listener() {
    auto exec = co_await this_coro::executor;
    tcp::acceptor acceptor(exec, {tcp::v4(), 12345});
    for (;;) {
        tcp::socket socket = co_await acceptor.async_accept(use_awaitable);
        // Her baglanti icin ayri bir coroutine baslat
        boost::asio::co_spawn(exec, echo(std::move(socket)),
                              boost::asio::detached);
    }
}

int main() {
    boost::asio::io_context io;
    boost::asio::co_spawn(io, listener(), boost::asio::detached);
    std::cout << "Coroutine yanki sunucusu 12345'te...\n";
    io.run();
    return 0;
}
\end{lstlisting}

\begin{lstlisting}[style=shstyle]
# C++20 ve coroutine destegi gerekir
g++ -std=c++20 -fcoroutines coro_echo.cpp -o coro_echo \
    -lboost_system -lpthread
\end{lstlisting}

\ipucu{Coroutine tabanlı Asio, modern C++ ağ programlamasının geleceğidir.
Yukarıdaki kod, önceki geri çağırmalı yankı sunucusuyla aynı işi yapar ama
akış çok daha nettir: \code{co\_await} her satırda ``burada asenkron bekle,
sonra devam et'' anlamına gelir. Hata yönetimi de normal \code{try/catch}
ile yapılır. Yeni Asio projelerine coroutine'lerle başlamanızı öneririm.}

% =====================================================================
\gsec{Boost.Atomic ve Boost.Lockfree}

\code{boost::lockfree}, kilit kullanmadan (lock-free) çalışan eşzamanlı veri
yapıları sunar: \code{queue}, \code{stack}, \code{spsc\_queue}. Yüksek
performanslı, düşük gecikmeli sistemlerde (finans, oyun, ses) kilitlerin
getirdiği bekleme ve öncelik tersine dönmesi (priority inversion)
sorunlarından kaçınmak için kullanılır.

\begin{lstlisting}
#include <boost/lockfree/queue.hpp>
#include <boost/thread.hpp>
#include <atomic>
#include <iostream>

int main() {
    // Kilitsiz kuyruk: coklu uretici/tuketici guvenli
    boost::lockfree::queue<int> queue(128);
    std::atomic<bool> done{false};
    std::atomic<long> total{0};

    // Uretici: 0..9999 arasi degerleri kuyruga koyar
    boost::thread producer([&]{
        for (int i = 0; i < 10000; ++i)
            while (!queue.push(i)) {}    // kuyruk doluysa tekrar dene
        done = true;
    });

    // Tuketici: kuyruktan cekip toplar
    boost::thread consumer([&]{
        int value;
        while (!done || !queue.empty()) {
            while (queue.pop(value))
                total += value;
        }
    });

    producer.join();
    consumer.join();
    std::cout << "Toplam: " << total << "\n";   // 49995000
    return 0;
}
\end{lstlisting}

\begin{lstlisting}[style=shstyle]
g++ -std=c++17 lockfree.cpp -o lockfree -lboost_thread -lboost_system -lpthread
\end{lstlisting}

\notk[SPSC Kuyruğu]{Tek üretici–tek tüketici (single producer, single
consumer) senaryonuz varsa \code{boost::lockfree::spsc\_queue} çok daha
hızlıdır; genel \code{queue}'nun CAS (compare-and-swap) döngülerine gerek
duymaz. Ses işleme ve gerçek zamanlı hatlarda tercih edilir.}

% #####################################################################
\gpart{5}{Metin Ayrıştırma ve Veri İşleme}
% #####################################################################

% =====================================================================
\gsec{Boost.Regex: Düzenli İfadeler}

Boost.Regex, \code{std::regex}'in atasıdır. Bugün \code{std::regex} standartta
olsa da, Boost.Regex genellikle \emph{daha hızlıdır}, Perl uyumluluğu daha
iyidir ve Unicode (ICU ile) desteği daha güçlüdür.

\begin{lstlisting}
#include <boost/regex.hpp>
#include <iostream>
#include <string>

int main() {
    std::string text = "iletisim: ali@example.com ve veli@test.org";

    // E-posta yakalayan duzenli ifade
    boost::regex email(R"((\w+)@(\w+\.\w+))");

    // Tum eslesmeleri gez
    boost::sregex_iterator it(text.begin(), text.end(), email);
    boost::sregex_iterator end;
    for (; it != end; ++it) {
        boost::smatch e = *it;
        std::cout << "Tam: " << e[0]        // ali@example.com
                  << " | kullanici: " << e[1]  // ali
                  << " | alan: " << e[2] << "\n";  // example.com
    }

    // Degistirme (replace) -- geri referans ile
    std::string masked = boost::regex_replace(
        text, email, "[gizli]@$2");
    std::cout << masked << "\n";
    return 0;
}
\end{lstlisting}

\begin{lstlisting}[style=shstyle]
g++ -std=c++17 regex_example.cpp -o regex_example -lboost_regex
\end{lstlisting}

% =====================================================================
\gsec{Boost.PropertyTree: JSON, XML, INI Ayrıştırma}

\code{boost::property\_tree}, hiyerarşik yapılandırma verilerini (JSON, XML,
INI, INFO) tek bir ağaç yapısında tutar ve bu formatlar arasında dönüşüm
yapar. Basit yapılandırma dosyaları için hızlı bir çözümdür.

\begin{lstlisting}
#include <boost/property_tree/ptree.hpp>
#include <boost/property_tree/json_parser.hpp>
#include <iostream>
#include <sstream>

namespace pt = boost::property_tree;

int main() {
    // JSON'dan bir property tree olustur
    std::string json = R"({
        "sunucu": {
            "host": "localhost",
            "port": 8080,
            "aktif": true
        },
        "kullanicilar": ["ali", "veli", "ayse"]
    })";

    pt::ptree tree;
    std::istringstream input(json);
    pt::read_json(input, tree);

    // Deger okuma (nokta ile yol belirtme)
    std::string host = tree.get<std::string>("sunucu.host");
    int port = tree.get<int>("sunucu.port");
    bool active = tree.get<bool>("sunucu.aktif");
    std::cout << host << ":" << port
              << " (aktif=" << std::boolalpha << active << ")\n";

    // Varsayilanli okuma (anahtar yoksa)
    int timeout = tree.get<int>("sunucu.timeout", 30);
    std::cout << "Zaman asimi: " << timeout << "\n";  // 30 (yok, varsayIlan)

    // Dizi (array) gezme
    std::cout << "Kullanicilar: ";
    for (const auto& [key, value] : tree.get_child("kullanicilar"))
        std::cout << value.get_value<std::string>() << " ";
    std::cout << "\n";

    // Degistir ve tekrar JSON yaz
    tree.put("sunucu.port", 9090);
    pt::write_json(std::cout, tree);
    return 0;
}
\end{lstlisting}

\uyari{\code{property\_tree} her değeri dahili olarak \emph{metin} olarak
saklar; bu yüzden JSON'un sayı, boolean ve dizi türlerini tam olarak
korumaz (örneğin diziler garip biçimde temsil edilir). Ciddi JSON işleri
için bir sonraki bölümdeki \code{boost::json}'u kullanın. PropertyTree'yi
basit INI/config dosyaları ve formatlar arası hızlı dönüşüm için saklayın.}

% =====================================================================
\gsec{Boost.JSON: Modern JSON İşleme}

Boost.JSON (Boost 1.75+), yüksek performanslı, standartlara uygun,
\emph{gerçek} bir JSON kütüphanesidir. PropertyTree'nin aksine JSON tiplerini
tam korur (sayı, string, bool, null, dizi, nesne) ve çok hızlıdır.

\begin{lstlisting}
#include <boost/json.hpp>
#include <iostream>
#include <string>

namespace json = boost::json;

int main() {
    // 1) Metinden ayristirma (parse)
    std::string text = R"({
        "isim": "Ada",
        "yas": 36,
        "diller": ["C++", "Python"],
        "aktif": true
    })";

    json::value v = json::parse(text);
    json::object& obj = v.as_object();

    std::cout << "Isim: " << obj["isim"].as_string() << "\n";
    std::cout << "Yas: "  << obj["yas"].as_int64() << "\n";
    std::cout << "Aktif: " << std::boolalpha << obj["aktif"].as_bool() << "\n";

    // Dizi gezme
    std::cout << "Diller: ";
    for (const auto& d : obj["diller"].as_array())
        std::cout << d.as_string() << " ";
    std::cout << "\n";

    // 2) Programatik olarak JSON insa etme
    json::object output;
    output["urun"]  = "Klavye";
    output["fiyat"] = 450.0;
    output["stok"]  = true;
    json::array tags = {"mekanik", "rgb"};
    output["etiketler"] = tags;

    // 3) Metne serilestirme (serialize)
    std::cout << "Uretilen JSON:\n" << json::serialize(output) << "\n";
    return 0;
}
\end{lstlisting}

\begin{lstlisting}[style=shstyle]
g++ -std=c++17 json_example.cpp -o json_example -lboost_json
\end{lstlisting}

\ipucu{Boost.JSON, \code{value\_to<T>} ve \code{value\_from<T>} ile kendi
C++ tiplerinizi doğrudan JSON'a/JSON'dan dönüştürebilir (kullanıcı tanımlı
dönüşümlerle). Ayrıca akış (streaming) ayrıştırıcısı sayesinde belleğe
sığmayan devasa JSON dosyalarını parça parça işleyebilir. Yeni projelerde
JSON için PropertyTree yerine Boost.JSON'u tercih edin.}

% =====================================================================
\gsec{Boost.Spirit: Ayrıştırıcı Üreteci}

Boost.Spirit, doğrudan C++ içinde, EBNF'ye benzer bir sözdizimiyle
ayrıştırıcı (parser) yazmanızı sağlayan güçlü bir kütüphanedir. Dilbilgisi
kurallarını C++ ifadeleri olarak yazarsınız; Spirit bunları derleme
zamanında bir ayrıştırıcıya dönüştürür. Küçük diller (DSL), yapılandırma
formatları ve protokol ayrıştırma için idealdir.

\begin{lstlisting}
#include <boost/spirit/home/x3.hpp>
#include <iostream>
#include <string>
#include <vector>

namespace x3 = boost::spirit::x3;

int main() {
    // Virgulle ayrilmis sayilari ayristir: "1.5, 2.0, 3.14"
    std::string input = "1.5, 2.0, 3.14, 42";
    std::vector<double> numbers;

    auto first = input.begin();
    auto end = input.end();

    // Dilbilgisi: double (% ile ayrilmis) -- bosluklari atla
    bool success = x3::phrase_parse(
        first, end,
        x3::double_ % ',',   // ',' ile ayrilmis double listesi
        x3::space,           // atlanacak bosluk kurali
        numbers              // sonuc buraya toplanIr (attribute)
    );

    if (success && first == end) {
        std::cout << "Ayristirildi (" << numbers.size() << " sayi): ";
        double sum = 0;
        for (double s : numbers) { std::cout << s << " "; sum += s; }
        std::cout << "\nToplam: " << sum << "\n";
    } else {
        std::cout << "Ayristirma basarisiz\n";
    }
    return 0;
}
\end{lstlisting}

\bilgi{Spirit X3 (yukarıdaki modern sürüm) header-only'dir ve derleme
zamanında çalışır. \code{x3::double\_ \% ','} ifadesi, ``virgülle ayrılmış
bir veya daha fazla double'' anlamına gelir. Spirit çok güçlüdür ama ağır
şablon kullanımı nedeniyle derleme süresi ve hata mesajları zorlayıcı
olabilir. Karmaşık dilbilgileri için harika, ama basit ayrıştırmalar için
\code{std::from\_chars} veya Tokenizer yeterli olabilir.}

% #####################################################################
\gpart{6}{Sistem, Matematik ve Graflar}
% #####################################################################

% =====================================================================
\gsec{Boost.Filesystem: Dosya Sistemi İşlemleri}

Boost.Filesystem, \code{std::filesystem}'in (C++17) atasıdır. Dosya ve dizin
işlemlerini platformdan bağımsız biçimde yapar: yol manipülasyonu, gezinme,
oluşturma/silme, boyut/zaman sorgulama.

\begin{lstlisting}
#include <boost/filesystem.hpp>
#include <iostream>

namespace fs = boost::filesystem;

int main() {
    fs::path p = "/tmp/boost_test";

    // Dizin olustur
    fs::create_directories(p / "alt/dizin");

    // Yol manipulasyonu
    fs::path file = p / "veri.txt";
    std::cout << "Ust dizin: " << file.parent_path() << "\n";
    std::cout << "Dosya adi: " << file.filename() << "\n";
    std::cout << "Uzanti: "    << file.extension() << "\n";

    // Varlik ve tur kontrolu
    std::cout << "Var mi? " << std::boolalpha << fs::exists(p) << "\n";
    std::cout << "Dizin mi? " << fs::is_directory(p) << "\n";

    // Dizini ozyinelemeli (recursive) gez
    std::cout << "\nIcerik:\n";
    for (const auto& entry : fs::recursive_directory_iterator(p))
        std::cout << "  " << entry.path()
                  << (fs::is_directory(entry) ? " [dizin]" : " [dosya]") << "\n";

    // Temizlik
    std::uintmax_t removed = fs::remove_all(p);
    std::cout << "\n" << removed << " oge silindi\n";
    return 0;
}
\end{lstlisting}

\begin{lstlisting}[style=shstyle]
g++ -std=c++17 fs_example.cpp -o fs_example -lboost_filesystem -lboost_system
\end{lstlisting}

\ipucu{C++17 hedefliyorsanız, çoğu durumda \code{std::filesystem} yeterlidir
ve bağımlılık gerektirmez. Boost.Filesystem'i, C++14/eski derleyici
desteklemeniz gerekiyorsa veya \code{std}'de henüz olmayan bazı ek
yardımcılara (örn. \code{unique\_path}, ek durum sorguları) ihtiyaç
duyuyorsanız tercih edin. API'ler neredeyse birebir aynı olduğundan geçiş
kolaydır.}

% =====================================================================
\gsec{Boost.DateTime ve Boost.Chrono}

Boost.DateTime, tarih ve saat işlemleri için zengin bir kütüphanedir:
takvim tarihleri, zaman aralıkları (duration), zaman dilimleri ve
biçimlendirme. \code{std::chrono}'nun (C++11/20) sunmadığı takvim ve zaman
dilimi özellikleri için hâlâ değerlidir.

\begin{lstlisting}
#include <boost/date_time/gregorian/gregorian.hpp>
#include <boost/date_time/posix_time/posix_time.hpp>
#include <iostream>

int main() {
    namespace greg = boost::gregorian;
    namespace pt   = boost::posix_time;

    // Tarih olusturma
    greg::date today(2026, greg::Jul, 19);
    greg::date newYear(2026, greg::Jan, 1);

    std::cout << "Bugun: " << today << "\n";
    std::cout << "Haftanin gunu: " << today.day_of_week() << "\n";
    std::cout << "Yilin kacinci gunu: " << today.day_of_year() << "\n";

    // Tarih aritmetigi
    greg::date_duration diff = today - newYear;
    std::cout << "Yilbasindan bu yana: " << diff.days() << " gun\n";

    greg::date future = today + greg::days(100);
    std::cout << "100 gun sonra: " << future << "\n";

    // Zaman (posix_time)
    pt::ptime now(today, pt::hours(14) + pt::minutes(30));
    std::cout << "Zaman damgasi: " << now << "\n";

    pt::time_duration duration = pt::hours(2) + pt::minutes(45) + pt::seconds(30);
    std::cout << "Sure: " << duration
              << " = " << duration.total_seconds() << " saniye\n";
    return 0;
}
\end{lstlisting}

\begin{lstlisting}[style=shstyle]
g++ -std=c++17 datetime_example.cpp -o datetime_example -lboost_date_time
\end{lstlisting}

% =====================================================================
\gsec{Boost.ProgramOptions: Komut Satırı Ayrıştırma}

Boost.ProgramOptions, komut satırı argümanlarını, ortam değişkenlerini ve
yapılandırma dosyalarını profesyonelce ayrıştırır. Otomatik yardım metni,
tip dönüşümü, varsayılan değerler ve doğrulama sağlar.

\begin{lstlisting}
#include <boost/program_options.hpp>
#include <iostream>
#include <string>

namespace po = boost::program_options;

int main(int argc, char* argv[]) {
    // Secenekleri tanImla
    po::options_description desc("Kullanim");
    desc.add_options()
        ("help,h", "Bu yardim metnini goster")
        ("isim,n", po::value<std::string>()->default_value("Dunya"),
                   "Selamlanacak isim")
        ("tekrar,t", po::value<int>()->default_value(1),
                     "Kac kez tekrarlansin")
        ("sesli,v", po::bool_switch(), "Ayrintili cikti");

    // Ayristir
    po::variables_map vm;
    try {
        po::store(po::parse_command_line(argc, argv, desc), vm);
        po::notify(vm);
    } catch (const po::error& e) {
        std::cerr << "Hata: " << e.what() << "\n";
        return 1;
    }

    // --help verildiyse yardimi goster ve cik
    if (vm.count("help")) {
        std::cout << desc << "\n";
        return 0;
    }

    // Degerleri kullan
    std::string name = vm["isim"].as<std::string>();
    int repeat = vm["tekrar"].as<int>();
    bool verbose = vm["sesli"].as<bool>();

    if (verbose) std::cout << "[ayrintili mod acik]\n";
    for (int i = 0; i < repeat; ++i)
        std::cout << "Merhaba, " << name << "!\n";
    return 0;
}
\end{lstlisting}

\begin{lstlisting}[style=shstyle]
g++ -std=c++17 opts.cpp -o opts -lboost_program_options
./opts --isim Ada --tekrar 3 -v
# [ayrintili mod acik]
# Merhaba, Ada!
# Merhaba, Ada!
# Merhaba, Ada!
./opts --help        # Otomatik uretilen yardim metnini gosterir
\end{lstlisting}

\ipucu{ProgramOptions ile aynı seçenekleri hem komut satırından hem de bir
\code{.cfg} dosyasından okuyabilirsiniz (\code{parse\_config\_file}).
Komut satırı, config dosyasını geçersiz kılacak şekilde katmanlayabilirsiniz;
bu, gerçek uygulamalarda standart bir yapılandırma stratejisidir.}

% =====================================================================
\gsec{Boost.Interprocess: Süreçler Arası İletişim}

Boost.Interprocess, süreçler arası iletişim (IPC) için paylaşımlı bellek,
bellek eşlemeli dosyalar, mesaj kuyrukları ve süreçler arası senkronizasyon
ilkelleri sunar. Aynı makinedeki farklı süreçlerin veri paylaşması gereken
yüksek performanslı senaryolarda kullanılır.

\begin{lstlisting}
#include <boost/interprocess/managed_shared_memory.hpp>
#include <iostream>

namespace ip = boost::interprocess;

int main() {
    // Onceki calismalardan kalmis olabilecek segmenti temizle
    ip::shared_memory_object::remove("BoostShm");

    // 64 KB'lik yonetilen paylasimli bellek segmenti olustur
    ip::managed_shared_memory segment(
        ip::create_only, "BoostShm", 65536);

    // Paylasimli bellekte bir int dizisi olustur ("Sayilar" adiyla)
    int* array = segment.construct<int>("Sayilar")[10](0);
    for (int i = 0; i < 10; ++i) array[i] = i * i;

    std::cout << "Paylasimli bellege yazildi.\n";

    // Baska bir surec ayni segmenti su sekilde acardi:
    //   managed_shared_memory seg(open_only, "BoostShm");
    //   auto res = seg.find<int>("Sayilar");
    //   int* dizi = res.first;  size_t n = res.second;

    // Ayni surec icinde bulup okuyalim (dogrulama)
    auto result = segment.find<int>("Sayilar");
    if (result.first) {
        std::cout << "Bulundu, " << result.second << " eleman: ";
        for (std::size_t i = 0; i < result.second; ++i)
            std::cout << result.first[i] << " ";
        std::cout << "\n";
    }

    ip::shared_memory_object::remove("BoostShm");   // temizle
    return 0;
}
\end{lstlisting}

\begin{lstlisting}[style=shstyle]
# Linux'ta paylasimli bellek icin rt kutuphanesi gerekebilir
g++ -std=c++17 ipc_example.cpp -o ipc_example -lrt -lpthread
\end{lstlisting}

\uyari{Paylaşımlı bellekte \emph{ham işaretçi} saklamayın; farklı süreçlerde
adresler farklı olacağından işaretçiler geçersiz olur. Bunun yerine
Boost.Interprocess'in \code{offset\_ptr} tipini ve paylaşımlı bellek uyumlu
ayırıcılarını (\code{allocator}) kullanın. Segmenti kullandıktan sonra
\code{remove} ile temizlemeyi unutmayın; aksi halde çekirdek nesnesi
sistemde kalır.}

% =====================================================================
\gsec{Boost.Multiprecision: Yüksek Hassasiyetli Aritmetik}

Boost.Multiprecision, \code{int} veya \code{double}'ın sınırlarını aşan
sayılarla çalışmanızı sağlar: keyfi hassasiyetli tam sayılar, ondalıklar
ve rasyonel sayılar. Kriptografi, bilimsel hesaplama ve tam sayı taşmasının
kabul edilemez olduğu durumlar için kritiktir.

\begin{lstlisting}
#include <boost/multiprecision/cpp_int.hpp>
#include <boost/multiprecision/cpp_dec_float.hpp>
#include <iostream>

namespace mp = boost::multiprecision;

// 100! -- normal int'te imkansiz (tasar), burada kesin
mp::cpp_int factorial(int n) {
    mp::cpp_int result = 1;
    for (int i = 2; i <= n; ++i)
        result *= i;
    return result;
}

int main() {
    // Keyfi hassasiyetli tam sayi
    std::cout << "100! = " << factorial(100) << "\n";

    // Cok buyuk uslu sayi
    mp::cpp_int big = 1;
    for (int i = 0; i < 200; ++i) big *= 2;     // 2^200
    std::cout << "2^200 = " << big << "\n";

    // 50 basamak hassasiyetli ondalik
    using dec50 = mp::number<mp::cpp_dec_float<50>>;
    dec50 pi = dec50("3.14159265358979323846264338327950288419716939937510");
    dec50 radius = 100;
    std::cout << "Alan = " << pi * radius * radius << "\n";

    // Cok hassasiyetli karekok
    dec50 two = 2;
    std::cout << "sqrt(2) = " << sqrt(two) << "\n";
    return 0;
}
\end{lstlisting}

\bilgi{\code{cpp\_int} header-only'dir ve harici bağımlılık gerektirmez.
Daha da yüksek performans için Boost, GMP (\code{gmp\_int}) ve MPFR
(\code{mpfr\_float}) arka uçlarını da destekler; bunlar ilgili sistem
kütüphanelerine bağlanmayı gerektirir ama devasa sayılarda çok daha
hızlıdır. Rasyonel sayılar için \code{cpp\_rational} kesin kesir aritmetiği
sunar.}

% =====================================================================
\gsec{Boost.Graph (BGL): Graf Algoritmaları}

Boost Graph Library (BGL), C++'ın en kapsamlı graf kütüphanesidir. Graf
gösterimlerini (komşuluk listesi/matrisi) ve üzerlerinde çalışan onlarca
klasik algoritmayı (BFS, DFS, Dijkstra, Bellman-Ford, Kruskal, Prim,
topolojik sıralama, güçlü bağlı bileşenler) sağlar. Genel programlama
(generic programming) tasarımı sayesinde kendi graf tipleriniz üzerinde
bile çalışabilir.

\gsub{Graf Oluşturma ve Gezinme}

\begin{lstlisting}
#include <boost/graph/adjacency_list.hpp>
#include <boost/graph/breadth_first_search.hpp>
#include <iostream>
#include <vector>

using namespace boost;

int main() {
    // Yonsuz graf: komsuluk listesi tabanlI
    using Graph = adjacency_list<vecS, vecS, undirectedS>;
    Graph g(6);   // 6 dugum (0..5)

    // Kenar ekle
    add_edge(0, 1, g);
    add_edge(0, 2, g);
    add_edge(1, 3, g);
    add_edge(2, 3, g);
    add_edge(3, 4, g);
    add_edge(4, 5, g);

    std::cout << "Dugum sayisi: " << num_vertices(g) << "\n";
    std::cout << "Kenar sayisi: " << num_edges(g) << "\n";

    // Tum kenarlari yazdir
    std::cout << "Kenarlar: ";
    auto [e_first, e_end] = edges(g);
    for (auto it = e_first; it != e_end; ++it)
        std::cout << source(*it, g) << "-" << target(*it, g) << " ";
    std::cout << "\n";

    // Bir dugumun komsularini gez
    std::cout << "0'in komsulari: ";
    auto [k_first, k_end] = adjacent_vertices(0, g);
    for (auto it = k_first; it != k_end; ++it)
        std::cout << *it << " ";
    std::cout << "\n";
    return 0;
}
\end{lstlisting}

\gsub{Dijkstra ile En Kısa Yol}

BGL'nin gerçek gücü hazır algoritmalarındadır. Aşağıda ağırlıklı bir grafta
en kısa yolları Dijkstra ile buluyoruz:

\begin{lstlisting}
#include <boost/graph/adjacency_list.hpp>
#include <boost/graph/dijkstra_shortest_paths.hpp>
#include <iostream>
#include <vector>

using namespace boost;

int main() {
    // Kenar agirligi tasiyan yonlu graf
    using Graph = adjacency_list<
        vecS, vecS, directedS,
        no_property,
        property<edge_weight_t, int>>;   // kenarlarda agirlik

    Graph g(5);
    // (kaynak, hedef, agirlik)
    add_edge(0, 1, 10, g);
    add_edge(0, 2, 5,  g);
    add_edge(2, 1, 3,  g);
    add_edge(1, 3, 1,  g);
    add_edge(2, 3, 9,  g);
    add_edge(3, 4, 4,  g);

    // Sonuc dizileri: her dugume uzaklik ve onceki dugum
    std::vector<int> distance(num_vertices(g));
    std::vector<Graph::vertex_descriptor> previous(num_vertices(g));

    auto start = vertex(0, g);

    // Dijkstra'yi calistir
    dijkstra_shortest_paths(g, start,
        predecessor_map(&previous[0]).distance_map(&distance[0]));

    // Sonuclari yazdir
    std::cout << "0'dan en kisa uzakliklar:\n";
    for (std::size_t i = 0; i < distance.size(); ++i)
        std::cout << "  Dugum " << i << ": " << distance[i] << "\n";

    // 0 -> 4 yolunu geri izle
    std::cout << "0'dan 4'e yol: ";
    std::vector<int> path;
    for (int v = 4; v != start; v = previous[v]) path.push_back(v);
    path.push_back(0);
    for (auto it = path.rbegin(); it != path.rend(); ++it)
        std::cout << *it << (it + 1 != path.rend() ? " -> " : "\n");
    return 0;
}
\end{lstlisting}

\begin{lstlisting}[style=shstyle]
# BGL cogunlukla header-only'dir; bu ornekler baglama gerektirmez
g++ -std=c++17 dijkstra.cpp -o dijkstra
./dijkstra
# 0'dan en kisa uzakliklar:
#   Dugum 0: 0
#   Dugum 1: 8   (0->2->1: 5+3)
#   Dugum 2: 5
#   Dugum 3: 9   (0->2->1->3: 5+3+1)
#   Dugum 4: 13
# 0'dan 4'e yol: 0 -> 2 -> 1 -> 3 -> 4
\end{lstlisting}

\bilgi{BGL, C++ genel programlamasının bir başyapıtıdır: algoritmalar graf
gösteriminden tamamen bağımsızdır ve ``property map'' soyutlaması sayesinde
düğüm/kenar verilerine esnek erişim sağlar. Öğrenme eğrisi diktir, ancak bir
kez kavradığınızda son derece güçlüdür. Karmaşık grafik hesaplamaları için
tekerleği yeniden icat etmek yerine BGL'yi değerlendirin.}

% =====================================================================
\gsec{Kapanış: Boost Ekosistemine Bakış}

Bu rehber Boost'un en çok kullanılan kütüphanelerini kapsadı, ancak Boost
bir buzdağıdır. Karşılaşabileceğiniz diğer önemli kütüphaneler:

\begin{center}
\begin{tabular}{@{}>{\raggedright\arraybackslash}p{3.6cm} >{\raggedright\arraybackslash}p{8.4cm}@{}}
\toprule
\textbf{Kütüphane} & \textbf{Amaç} \\
\midrule
\code{Boost.Hana} & Modern (C++14+) metaprogramlama; MPL'nin halefi. \\
\code{Boost.Mp11} & Hafif, pratik tip-listesi metaprogramlama. \\
\code{Boost.Geometry} & Uzamsal algoritmalar, R-tree, coğrafi hesap. \\
\code{Boost.Beast} & Asio üzerine HTTP ve WebSocket. \\
\code{Boost.Log} & Esnek, ölçeklenebilir günlükleme (logging). \\
\code{Boost.Serialization} & Nesneleri diske/ağa serileştirme. \\
\code{Boost.Fiber} & Kullanıcı-alanı iş parçacıkları (M:N zamanlama). \\
\code{Boost.Signals2} & İş parçacığı güvenli sinyal/yuva (observer). \\
\code{Boost.Coroutine2} & Yığın tabanlı (stackful) coroutine'ler. \\
\code{Boost.Test} & Kapsamlı birim test çerçevesi. \\
\code{Boost.Math} & Özel fonksiyonlar, istatistik, dağılımlar. \\
\code{Boost.uBLAS} & Temel doğrusal cebir (matris/vektör). \\
\bottomrule
\end{tabular}
\end{center}

\ipucu{\textbf{Öğrenme stratejisi:} Boost'un tamamını öğrenmeye çalışmayın;
bu ne mümkün ne de gerekli. Bunun yerine, bir problemle karşılaştığınızda
``Bunun için bir Boost kütüphanesi var mı?'' diye sorun. Genellikle vardır
ve iyi test edilmiştir. Her kütüphanenin resmi dokümantasyonu, çalışan
örneklerle doludur ve başlangıç noktanız olmalıdır: \texttt{boost.org/doc/libs}.}

\bilgi{\textbf{Modern C++ ve Boost dengesi:} C++ standardı her sürümde
büyüdükçe (C++20 concepts, ranges, coroutines; C++23 \code{flat\_map},
\code{expected}), Boost'un bazı kütüphaneleri gereksizleşiyor. Yine de Asio,
Spirit, Graph, Multiprecision, Geometry, Beast ve Interprocess gibi
kütüphaneler standartta karşılığı olmayan, vazgeçilmez araçlar olmaya devam
ediyor. Doğru yaklaşım: standartta varsa \code{std}'yi, yoksa Boost'u
kullanın.}

\vspace{1cm}
\begin{center}
{\color{ortaGri}\rule{0.5\linewidth}{0.8pt}}\\[6pt]
{\color{anaMavi}\large\bfseries Rehberin Sonu}\\[4pt]
{\color{ortaGri}\small Boost 1.85+ \textbullet\ C++17/20 \textbullet\ İyi kodlamalar!}
\end{center}

\end{document}
